%=== Câu 5 (ID: None) ===%
\begin{ex}
Trong các câu sau có bao nhiêu câu là mệnh đề
\begin{enumEX}{2}
\item Ngôi nhà đẹp quá!
\item Năm 1900 không phải là năm nhuận.
\item Số 9 là số nguyên tố.
\item $x^3+1=0$.
\end{enumEX}
\choice
{$1$}
{$3$}
{\True $2$}
{$4$}
\loigiai{
\lq\lq Ngôi nhà đẹp quá!\rq\rqlà câu cảm thán, không phải là mệnh đề.\\
\lq\lq $x^3+1=0$\rq\rqlà mệnh đề chứa biến, không phải là mệnh đề.}
\end{ex}

%=== Câu 9 (ID: None) ===%
\begin{ex}%%[0D1N1-1]%[Dự án đề kiểm tra Toán 10 HKI NH23-24- Nguyễn Hoàng Anh]%[THPT Nguyễn Chí Thanh - Lâm Đồng]
Trong các câu sau, câu nào là mệnh đề?
\choice
{\True Số $5$ là số nguyên tố}
{Lớp 10A1 có bao nhiêu học sinh?}
{$x+2=5$}
{$x^2 {<} 4$}
\loigiai{ Câu "Số $5$ là số nguyên tố" là một khẳng định đúng nên nó là một mệnh đề.
}
\end{ex}

%=== Câu 12 (ID: None) ===%
\begin{ex}
Có bao nhiêu phát biểu là mệnh đề?
\begin{itemize}
\item Hà Nội là Thủ đô của Việt Nam.
\item Số 16 là một số chính phương.
\item $x=1$ có phải là nghiệm của phương trình $x^2-1=0$ không?
\item Phương trình $3x^2-5x+2=0$ có nghiệm nguyên.
\item $5{<}7-3$.
\end{itemize}
\choice
{$2$}
{\True $4$}
{$3$}
{$5$}
\loigiai{Có $4$ phát biểu là mệnh đề đó là
\begin{itemize}
\item Hà Nội là Thủ đô của Việt Nam.
\item Số 16 là một số chính phương.
\item Phương trình $3x^2-5x+2=0$ có nghiệm nguyên.
\item $5{<}7-3$.
\end{itemize}
}
\end{ex}

%=== Câu 27 (ID: None) ===%
\begin{ex}
Trong các câu sau, câu nào là mệnh đề?
\def\dotEX{}% Thay đổi dấu chấm khi kết thúc mỗi p/a trả lời
\choice
{Đề thi môn Toán khó quá!}
{Bạn có đi học không?}
{Mùa thu Hà Nội đẹp quá!}
{\True Hà Nội là thủ đô của Việt Nam.}
\loigiai{
Phát biểu \lq\lq Hà Nội là thủ đô của Việt Nam\rq\rq\,  là một mệnh đề.
}
\end{ex}

%=== Câu 59 (ID: None) ===%
\begin{ex}
Mệnh đề ``$\exists x \in \mathbb{R}\colon x^2=9$'' khẳng định rằng
\choice
{Nếu $x$ là số thực thì $x^2=9$}
{\True Có ít nhất một số thực mà bình phương của nó bằng $9$}
{Bình phương của một số thực bằng $9$}
{Chỉ có một số thực bình phương bằng $9$}
\loigiai{
Mệnh đề ``$\exists x \in \mathbb{R}\colon x^2=9$'' khẳng định rằng: Có ít nhất một số thực mà bình phương của nó bằng $9$.
}
\end{ex}

%=== Câu 60 (ID: None) ===%
\begin{ex}
Cho mệnh đề chứa biến $P(x)\colon \lq\lq x^2{<}3x$ với $x$ là số thực. Mệnh đề nào sau đây đúng?
\choice
{$P(0)$}
{$P(5)$}
{$P(-1)$}
{\True $P(2)$}
\loigiai{
Thay các giá trị $2, 0, 5, -1$ ta được $2^2{<}3\cdot 2$ nên $P(2)$ đúng.
}
\end{ex}

%=== Câu 62 (ID: None) ===%
\begin{ex}
Ký hiệu nào sau đây dùng để viết đúng mệnh đề ``$\sqrt{5}$ không phải là số hữu tỉ''?
\choice
{\True $\sqrt{5} \notin \mathbb{Q}$}
{$\sqrt{5} \notin \mathbb{R}$}
{$\sqrt{5} \neq \mathbb{R}$}
{$\sqrt{5} \not \subset \mathbb{Q}$}
\loigiai{
Mệnh đề ``$\sqrt{5}$ không phải là số hữu tỉ'' được kí hiệu là ``$\sqrt{5} \notin \mathbb{Q}$''.
}
\end{ex}

%=== Câu 63 (ID: None) ===%
\begin{ex}
Cho mệnh đề $B\colon$ \lq\lq$\exists x \in \mathbb{R}, x^2 = 2$\rq\rq. Phát biểu bằng lời của mệnh đề $B$ là
\choice
{\True Có ít nhất một số thực mà bình phương của nó bằng $2$}
{Bình phương của mỗi số thực đều bằng $2$}
{Nếu $x$ là số thực thì $x^2$ bằng $2$}
{Chỉ có một số thực mà bình phương của nó bằng $2$}
\loigiai{
Phát biểu bằng lời của mệnh đề $B\colon$ \lq\lq$\exists x \in \mathbb{R}, x^2 = 2$\rq\rq\, là \lq\lq Có ít nhất một số thực mà bình phương của nó bằng $2$\rq\rq.
}
\end{ex}

%=== Câu 76 (ID: None) ===%
\begin{ex}
Cho mệnh đề chứa biến $P(n):``n^2-1$ chia hết cho $4$ với $n$ là số nguyên. Xét các mệnh đề $P(5)$ và $P(2)$ đúng hay sai?
\choice
{$P(5)$ sai và $P(2)$ sai}
{$P(5)$ đúng và $P(2)$ đúng}
{$P(5)$ sai và $P(2)$ đúng}
{\True $P(5)$ đúng và $P(2)$ sai}
\end{ex}

%=== Câu 91 (ID: None) ===%
\begin{ex}
Mệnh đề là:
\choice
{Một khẳng định luôn đúng}
{Câu cảm thán}
{Câu nghi vấn hoặc câu cầu khiến}
{\True Một khẳng định chỉ có thể đúng hoặc sai}
\loigiai{
Mệnh đề là một khẳng định chỉ có thể đúng hoặc sai.
}
\end{ex}

%=== Câu 100 (ID: None) ===%
\begin{ex}
Xét tính đúng sai của các khẳng định sau.
\choiceTF
{Phát biểu "Ở đây đẹp quá!" là một mệnh đề}
{\True Phát biểu "Phương trình $x^2-3 x+1=0$ vô nghiệm." là một mệnh đề}
{\True Phát biểu "Số $16$ không phải là số nguyên tố." là một mệnh đề}
{Phát biểu "Số $\pi$ có lớn hơn $3$ hay không?" là một mệnh đề}
\loigiai{
\begin{itemchoice}
\itemch sai.
\itemch đúng.
\itemch đúng.
\itemch đúng.
\end{itemchoice}
}
\end{ex}

%=== Câu 103 ===%
\begin{ex}
Trong các mệnh đề sau, có bao nhiêu mệnh đề toán học?  $\bullet$ $A\colon$ '' Hai đường thẳng song song không có điểm chung ''.  $\bullet$ $B\colon$ '' Tích hai số thực trái dấu là một số thực âm ''.  $\bullet$ $C\colon$ '' Hàm số $y = x^2$ có đồ thị là một parabol ''.  $\bullet$ $D\colon$ '' Mặt Trời quay quanh Trái Đất ''.  $\bullet$ $E\colon$ '' Mọi số tự nhiên đều là dương ''.
\shortans{4}
\loigiai{$\bullet$ Phát biểu '' Hai đường thẳng song song không có điểm chung '' là một mệnh đề toán học.  $\bullet$ Phát biểu '' Tích hai số thực trái dấu là một số thực âm '' là một mệnh đề toán học.  $\bullet$ Phát biểu '' Hàm số $y = x^2$ có đồ thị là một parabol '' là một mệnh đề toán học.  $\bullet$ Phát biểu '' Mặt Trời quay quanh Trái Đất '' không là một mệnh đề toán học (vì không liên quan đến sự kiện Toán học nào).  $\bullet$ Phát biểu '' Mọi số tự nhiên đều là dương '' là một mệnh đề toán học.}
\end{ex}

%=== Câu 105 ===%
\begin{ex}
Có bao nhiêu câu sau đây là mệnh đề ?  $\bullet$ $A\colon$ '' Ôi! Trời lạnh thật đấy! ''.  $\bullet$ $B\colon$ '' Trái Đất hình lập phương ''.  $\bullet$ $C\colon$ '' Hàm số $y = x^2$ có đồ thị là một parabol ''.  $\bullet$ $D\colon$ '' Pi là một số vô tỷ ''.  $\bullet$ $E\colon$ '' Mọi số tự nhiên đều là dương ''.
\shortans{4}
\loigiai{$\bullet$ Phát biểu '' Ôi! Trời lạnh thật đấy! '' không phải là một mệnh đề vì không thể xác định đúng hoặc sai.  $\bullet$ Phát biểu '' Trái Đất hình lập phương '' là một mệnh đề vì có thể xác định đúng hoặc sai.  $\bullet$ Phát biểu '' Hàm số $y = x^2$ có đồ thị là một parabol '' là một mệnh đề vì có thể xác định đúng hoặc sai.  $\bullet$ Phát biểu '' Pi là một số vô tỷ '' là một mệnh đề vì có thể xác định đúng hoặc sai.  $\bullet$ Phát biểu '' Mọi số tự nhiên đều là dương '' là một mệnh đề vì có thể xác định đúng hoặc sai.}
\end{ex}
